% !TeX root = ../nc-quadric-avatar.tex
% 2026-11-22: relative Azumaya vertex classes and finite directed Morita model
\section{Brauer重みを持つ四頂点Morita表示}
\label{sec:qa-twisted-morita-weights}
前節のgerbe上の幾何avatarの四項Endは、慣性の重みを
記録してAzumaya頂点・双加群矢により係数空間へ降下する。
ここでは有限EndのMorita表示のみを扱い、
相対AS完成およびpointed \qgr の同値は主張しない。
gerbeの重み分解およびtwisted sheafのMorita理論は
\cite[Theorems~4.7 and 5.4]{BerghSchnurer2021}
および\cite[Theorem~6.2 and Remark~6.3]{BerghSchnurer2021}
の一般論を使う。

\begin{lemma}[conic pencilのdeterminant補正]
\label{lem:qa-morita-determinant-normalization}
前節の$U_2$と$\mathscr G=\mathscr G_c\times_{U_2}\mathscr G_d$を取る。
$q_i:Q_i\to\mathscr G$を普遍smooth conic、
$N_i$を普遍相対幾何次数一束とし、
$V_i=q_{i*}N_i$（$i=c,d$）と置く。
$V_i$は階数二、全基底変更可換で、
一次元商の規約で$Q_i\simeq\PP_{\mathscr G}(V_i)$となる。
conic超平面束を$\OO_{Q_i}(1)$とすれば
\begin{equation}
 N_i^{\otimes2}\simeq
 \OO_{Q_i}(1)\otimes q_i^*\det V_i .
 \label{eq:qa-morita-determinant-square}
\end{equation}
$E=Q_c\times_{\PP^1}Q_d$、
$p_i:E\to Q_i$、$\pi:E\to\mathscr G$、
$F=p_c^*N_c$、$G=p_d^*N_d$、
$H=\OO_E(1)$と書くと
\begin{equation}
 F^2\simeq H\pi^*\det V_c,\quad
 G^2\simeq H\pi^*\det V_d,\quad
 (GF^{-1})^2\simeq\pi^*(\det V_d(\det V_c)^{-1}).
 \label{eq:qa-morita-fg-square}
\end{equation}
従って$GF^{-1}$が相対Picardの二捩れsectionであることと、
gerbe上でその二乗線束が実際に自明であることは異なる。
\end{lemma}
\begin{proof}
各幾何conicは$\PP^1$であり、$N_i$の相対次数は一だから
$V_i$は階数二で$R^1q_{i*}N_i=0$、
相対評価は全射である。評価が定める
$Q_i\to\PP(V_i)$は各fiberで次数一の完全線形系なので同型。
商規約の相対Euler列は
$\omega_{Q_i/\mathscr G}\simeq N_i^{-2}\otimes q_i^*\det V_i$
を与える。一方、二次Cartier超曲面
$Q_i\subset\PP^2_{U_2}$の随伴から
$\omega_{Q_i/\mathscr G}\simeq\OO_{Q_i}(-1)$である。
これで最初の式が得られる。
両conicの超平面束は、共通の$\PP^1_{[u:v]}$の
$\OO(1)$の引戻しだから、$E$上で同じ$H$になる。
二式の引戻しと商で残りの式が従う。
$\det V_c,\det V_d$の慣性重みは$(2,0),(0,2)$であり、
最後の右辺の自明化を仮定してはならない。
\end{proof}

\begin{theorem}[四頂点のAzumaya双加群による有限Morita表示]
\label{thm:qa-twisted-four-vertex-morita}
前補題の記号で、gerbe上の四線束を
\begin{equation}
 (L_0,L_1,L_2,L_3)=(\OO_E,F,H,HG)
 \label{eq:qa-morita-normalized-four}
\end{equation}
と正規化する。
$T_{ij}=\pi_*(L_i^{-1}L_j)$（$i\le j$）は
有限局所自由で、任意の基底変更と可換する。
非対角成分の階数と慣性$\Gm^2$の指標は
\begin{equation}
\begin{array}{c|rrrrrr}
(i,j)&01&02&03&12&13&23\\\hline
\operatorname{rk}T_{ij}&2&4&6&2&4&2\\
\operatorname{wt}T_{ij}&(1,0)&(0,0)&(0,1)&(-1,0)&(-1,1)&(0,1)
\end{array}
\label{eq:qa-morita-hom-weights}
\end{equation}
であり、$T_{ii}=\OO_{\mathscr G}$である。
四項の重みは
$w_0=(0,0)$、$w_1=(1,0)$、$w_2=(0,0)$、$w_3=(0,1)$。
$W_0=W_2=\OO_{\mathscr G}$、
$W_1=V_c$、$W_3=V_d$とする。
重み零の降下を$\operatorname{Desc}$で記し、
\[
 A_i=\operatorname{Desc}\End(W_i),\qquad
 {\mathsf M}_{ij}
 =\operatorname{Desc}(T_{ij}\otimes W_i\otimes W_j^\vee)
 \quad(i\le j)
\]
と定義する。$A_i$は$U_2$上のAzumaya代数であり
\begin{equation}
 \operatorname{rk}(A_0,A_1,A_2,A_3)=(1,4,1,4),\qquad
 ([A_0],[A_1],[A_2],[A_3])=(0,\alpha_c,0,\alpha_d).
 \label{eq:qa-morita-vertex-classes}
\end{equation}
（oppositeの規約による符号は二捩れのため影響しない。）
$\mathsf M_{ij}$は有限局所自由な$(A_i,A_j)$双加群、
$\mathsf M_{ii}=A_i$であり、
\begin{equation}
 {\mathsf M}_{ij}\otimes_{A_j}{\mathsf M}_{jk}
 \longrightarrow{\mathsf M}_{ik}
 \qquad(i\le j\le k)
 \label{eq:qa-morita-bimodule-product}
\end{equation}
は結合的で単位を持つ。
この四頂点Azumaya双加群有向圏は
$U_2$の任意の基底変更と可換する。
分岐根を選ぶ有限étale分裂被覆の上で
$V_i$と$\det V_c$をさらに局所自明化すると、
命題~\ref{prop:qa-point-relative-split-root-end}の
spinor有限End代数の通常のMorita増幅に一致する。
\emph{相対AS完成と相対pointed $\qgr$の同値は結論しない。}
\end{theorem}
\begin{proof}
$\pi:E\to\mathscr G$は$(2,2)$ Cartier族である。
$i<j$なら$L_jL_i^{-1}$の各既約成分上の次数は正である。
相対$(2,2)$分解とcohomology-and-base-changeによって
$R^1\pi_*(L_jL_i^{-1})=0$となり、
$T_{ij}$は指定階数の有限局所自由束で基底変更と可換する。
対角は$\pi_*\OO_E=\OO_{\mathscr G}$である
（$R^1\pi_*\OO_E$の消滅は用いない）。

gerbeの慣性$(\lambda_c,\lambda_d)$は
$N_c,N_d$を各々の文字でスカラー倍する。
$F,G$の重みは$(1,0),(0,1)$、
二次形式の係数に由来する$H$の重みは$(0,0)$である。
$T_{ij}$の重みは$w_j-w_i$となり、
表の全項が従う。

$W_i$の重みは$w_i$である。
$\End(W_i)$は重み零、
$T_{ij}\otimes W_i\otimes W_j^\vee$の重みは
$(w_j-w_i)+w_i-w_j=0$である。
対角化可能群をbandとするgerbeの重み零sheafは
基底へ一意に降下する
\cite[Theorem~4.7]{BerghSchnurer2021}。
そのため$A_i$と$\mathsf M_{ij}$が存在する。
$V_c,V_d$はそれぞれのBrauer類のrank-two twisted bundleなので、
$\End(V_c),\End(V_d)$の降下は階数四Azumayaである。

$W_i$への左$\End(W_i)$作用と
$W_j^\vee$への右$\End(W_j)$作用で
$\mathsf M_{ij}$を双加群にする。
合成には線束の乗法$T_{ij}T_{jk}\to T_{ik}$と
\[
 (W_i\otimes W_j^\vee)
 \otimes_{\End W_j}(W_j\otimes W_k^\vee)
 \simeq W_i\otimes W_k^\vee
\]
というMorita evaluationを用いる。
結合律と単位はこの二つの構成のそれから従う。
重み零なので降下後も各等式は成り立つ。

分岐根選択の被覆上では両conicが$\PP^1$-bundleに分裂し、
各$V_i$を局所自明化できる。
すると$A_i$は行列代数、$\mathsf M_{ij}$は
$T_{ij}$のMorita増幅になる。
式~\eqref{eq:qa-morita-fg-square}で
$\det V_c$を自明化すると四項は
$(\OO_E,F,F^2,F^2G)$に一致する。
命題~\ref{prop:qa-point-relative-split-root-end}を適用して
局所spinor有限EndとのMorita同型が得られる。
一般の全整数AS完成の降下は使っていない。
\end{proof}

\begin{corollary}[ruling交換の頂点固定Morita降下は不可能]
\label{cor:qa-morita-ruling-exchange-phase}
$U_2$の二つの二次形式の交換$\epsilon$は
上記の頂点Brauer類を
\[
 (0,\alpha_c,0,\alpha_d)
   \longmapsto(0,\alpha_d,0,\alpha_c)
\]
と交換する。
一般点では$\alpha_c,\alpha_d$が独立なので、
この交換を全四頂点\emph{固定}のMorita自己同値として
降下させることは不可能である。
特にO-pointingを固定した相対比較でも、
残りのcollectionの位相変更・mutationを省略して
頂点固定同型を主張してはならない。
適切なmutationを伴う同値の存在・不存在は
この結論に含まない。
\end{corollary}
\begin{proof}
頂点固定のMorita同値が存在すれば、
特に一番頂点のAzumaya代数$A_1$と
$\epsilon^*A_1$はMorita同値となる。
そのBrauer類はそれぞれ$\alpha_c,\alpha_d$なので、
$\alpha_c=\alpha_d$が必要である。
しかし命題~\ref{prop:qa-universal-conic-brauer-obstruction}は
一般係数体上で$\alpha_c-\alpha_d\ne0$を証明した。
この矛盾は頂点固定の同値だけを排除し、
異なる四項の導来mutationには適用されない。
\end{proof}
